% ---------------------------------------------------------------------
\chapter{Fundamentos: el electrocardiograma y el síndrome de Brugada}\label{cap:fundamentos}

Este capítulo resume los conceptos clínicos necesarios para entender las decisiones de los capítulos siguientes: qué mide un ECG, en qué consiste el síndrome de Brugada y por qué su patrón es difícil de detectar.

\section{El electrocardiograma}

El electrocardiograma (ECG) registra en la superficie del cuerpo la actividad eléctrica que acompaña a la contracción del corazón. Cada latido produce una secuencia de ondas (\cref{fig:latido}):

\begin{itemize}
  \item la \textbf{onda P}, despolarización de las aurículas;
  \item el \textbf{complejo QRS}, despolarización de los ventrículos, con una duración normal por debajo de unos \SI{120}{\milli\second};
  \item el \textbf{punto J}, que marca el final del QRS y el comienzo del segmento ST;
  \item el \textbf{segmento ST}, una meseta normalmente isoeléctrica durante la cual los ventrículos permanecen despolarizados;
  \item la \textbf{onda T}, repolarización de los ventrículos.
\end{itemize}

El ECG estándar tiene 12 derivaciones, cada una de las cuales observa el mismo fenómeno eléctrico desde un ángulo distinto: seis de miembros (I, II, III, aVR, aVL, aVF) en el plano frontal y seis precordiales (V1--V6) en el plano horizontal. Las \textbf{precordiales derechas V1 y V2} se sitúan sobre el tracto de salida del ventrículo derecho, que es la región donde se origina el síndrome de Brugada. Por eso son las derivaciones en las que aparece su patrón y en las que se centra este trabajo.

\begin{figure}[htbp]
\centering
\begin{tikzpicture}[x=1.15cm, y=1.9cm, font=\sffamily\footnotesize]
  % rejilla
  \foreach \x in {0,0.4,...,10.01} \draw[acentoclaro, line width=0.3pt] (\x,-0.55) -- (\x,1.35);
  \foreach \y in {-0.5,-0.25,...,1.3} \draw[acentoclaro, line width=0.3pt] (0,\y) -- (10,\y);
  \draw[gristexto!50, dashed, line width=0.4pt] (0,0) -- (10,0);
  % latido normal
  \draw[acento, line width=1.3pt, line join=round]
    (0,0) -- (0.6,0) .. controls (0.9,0.2) and (1.3,0.2) .. (1.6,0) -- (2.6,0)
    -- (2.8,-0.12) -- (3.15,1.15) -- (3.5,-0.38) -- (3.72,0) -- (4.7,0.02)
    .. controls (5.3,0.05) and (5.6,0.48) .. (6.1,0.42) .. controls (6.5,0.38) and (6.8,0) .. (7.3,0) -- (10,0);
  % etiquetas de ondas
  \node[above] at (1.1,0.16) {P};
  \node[above] at (3.15,1.15) {R};
  \node[below left] at (2.85,-0.12) {Q};
  \node[below right] at (3.45,-0.38) {S};
  \node[above] at (6.1,0.45) {T};
  \fill[brugada] (3.72,0) circle (1.6pt) node[below right, text=brugada] {J};
  % llaves
  \draw[decorate, decoration={brace, amplitude=4pt, mirror}, gristexto] (2.6,-0.62) -- (3.72,-0.62) node[midway, below=4pt] {QRS};
  \draw[decorate, decoration={brace, amplitude=4pt, mirror}, gristexto] (3.72,-0.62) -- (5.0,-0.62) node[midway, below=4pt] {segmento ST};
  \draw[decorate, decoration={brace, amplitude=4pt}, gristexto] (0.6,0.42) -- (2.6,0.42) node[midway, above=4pt] {intervalo PR};
  % ventana de análisis
  \draw[brugada!70, line width=0.5pt, dashed] (3.72,-0.5) -- (3.72,1.3);
  \draw[brugada!70, line width=0.5pt, dashed] (4.62,-0.5) -- (4.62,1.3);
  \node[text=brugada, font=\sffamily\scriptsize, align=center] at (4.17,1.12) {J a J+80\\ ms};
\end{tikzpicture}
\caption[Anatomía de un latido en el ECG]{Anatomía esquemática de un latido. Cada cuadrícula horizontal representa \SI{40}{\milli\second}. El punto J (en rojo) separa la despolarización ventricular (QRS) de la repolarización (segmento ST y onda T). La ventana discontinua, de J a J+\SI{80}{\milli\second}, es la región donde se manifiesta el patrón de Brugada y donde se miden muchas de las características del \cref{cap:caracteristicas}.}
\label{fig:latido}
\end{figure}

\section{El síndrome de Brugada}

\subsection{Mecanismo, epidemiología y presentación clínica}

El síndrome de Brugada (SBr) es una \textbf{canalopatía}: una enfermedad de los canales iónicos de la membrana de las células cardíacas, sin alteración estructural visible del corazón. En torno a una quinta parte de los pacientes presenta variantes con pérdida de función en el gen \emph{SCN5A}, que codifica el canal de sodio cardíaco. La reducción de la corriente de sodio se expresa sobre todo en el epicardio del tracto de salida del ventrículo derecho \cite{brugada2018jacc,antzelevitch2016}. Coexisten dos explicaciones, no excluyentes, de cómo se genera el patrón: un gradiente de repolarización entre el epicardio y el endocardio de esa región (hipótesis de repolarización) y un retraso de conducción localizado (hipótesis de despolarización).

Su prevalencia se estima entre 1 de cada \num{2000} y 1 de cada \num{5000} personas, con valores más altos en el sudeste asiático. Es entre ocho y diez veces más frecuente en varones, y suele diagnosticarse en torno a los 40~años \cite{brugada2018jacc}. Muchos pacientes están asintomáticos. Otros debutan con síncope o con fibrilación ventricular, típicamente en reposo, durante el sueño o con fiebre.

\subsection{Patrones electrocardiográficos}

Los consensos internacionales describen dos patrones en V1--V3 \cite{wilde2002,bayesdeluna2012,priori2013} (\cref{fig:patrones}):

\begin{itemize}
  \item \textbf{Tipo 1 o \eng{coved} (convexo)}: elevación del punto J de al menos \SI{2}{\milli\metre} (\SI{0.2}{\milli\volt}), seguida de un segmento ST descendente y convexo que termina en una \textbf{onda T negativa} y simétrica. Es el único patrón diagnóstico.
  \item \textbf{Tipo 2 o \eng{saddleback} (en silla de montar)}: onda r' elevada, seguida de un segmento ST cóncavo que se mantiene elevado al menos \SI{0.5}{\milli\metre} y de una onda T positiva o aplanada. Es sugestivo, pero no diagnóstico. El antiguo tipo 3 se integró en él en 2012 \cite{bayesdeluna2012}.
\end{itemize}

Para distinguir el tipo 2 de variantes benignas con onda r' (bloqueo incompleto de rama derecha, deportistas) se han propuesto criterios cuantitativos sobre la forma del complejo: el ángulo $\beta$ entre las dos ramas de la r' \cite{chevallier2011} y la anchura de la base del triángulo que forma la r' \cite{serra2014}. Para distinguir el patrón tipo 1 de la repolarización precoz de los deportistas se propuso el \textbf{índice de Corrado}, el cociente entre la elevación del ST en el punto J y a \SI{80}{\milli\second} ($\mathrm{ST_J}/\mathrm{ST_{80}}$): un valor mayor que 1 indica un ST descendente, propio del patrón de Brugada \cite{zorzi2015}. Todos estos criterios miden la \emph{forma} de la transición entre el QRS y el ST, no su altura. Los datos de este trabajo confirman empíricamente esa idea (\cref{cap:caracteristicas}).

\begin{figure}[htbp]
\centering
\begin{tikzpicture}[x=0.95cm, y=1.55cm, font=\sffamily\footnotesize]
  \foreach \ox/\tit in {0/{Normal (V1)}, 5.2/{Tipo 1 (\eng{coved})}, 10.4/{Tipo 2 (\eng{saddleback})}} {
    \begin{scope}[shift={(\ox,0)}]
      \foreach \x in {0,0.4,...,4.41} \draw[acentoclaro, line width=0.3pt] (\x,-1.1) -- (\x,1.3);
      \foreach \y in {-1.0,-0.75,...,1.3} \draw[acentoclaro, line width=0.3pt] (0,\y) -- (4.4,\y);
      \draw[gristexto!50, dashed, line width=0.4pt] (0,0) -- (4.4,0);
      \node[anchor=south west, font=\sffamily\small\bfseries, text=acento] at (0,1.32) {\tit};
    \end{scope}}
  % normal en V1: rS, ST isoeléctrico, T positiva baja
  \draw[normal, line width=1.3pt, line join=round] (0,0) -- (0.9,0) -- (1.05,0.25) -- (1.25,-0.95) -- (1.45,0) -- (2.1,0.03)
    .. controls (2.5,0.08) and (2.7,0.3) .. (3.0,0.28) .. controls (3.3,0.26) and (3.5,0) .. (3.8,0) -- (4.4,0);
  % tipo 1: r' alta que desciende en rampa convexa hasta T negativa
  \begin{scope}[shift={(5.2,0)}]
    \draw[brugada, line width=1.3pt, line join=round] (0,0) -- (0.9,0) -- (1.05,0.2) -- (1.22,-0.7) -- (1.42,0.95)
      .. controls (1.9,0.92) and (2.35,0.65) .. (2.65,0.05) .. controls (2.85,-0.35) and (3.1,-0.5) .. (3.35,-0.3)
      .. controls (3.55,-0.12) and (3.7,0) .. (3.95,0) -- (4.4,0);
    \draw[gristexto, -{Latex[length=1.6mm]}, line width=0.5pt] (2.75,1.08) node[right, font=\sffamily\scriptsize, text=gristexto] {ST descendente} -- (2.12,0.84);
    \draw[gristexto, -{Latex[length=1.6mm]}, line width=0.5pt] (3.6,-0.85) node[below, font=\sffamily\scriptsize, text=gristexto] {T negativa} -- (3.2,-0.47);
    \draw[brugada, line width=0.5pt] (1.42,0) -- (1.42,0.95); \node[left, font=\sffamily\scriptsize, text=brugada, fill=white, inner sep=1pt] at (1.36,0.62) {$\geq$\,2 mm};
  \end{scope}
  % tipo 2: r' y ST en silla de montar, T positiva
  \begin{scope}[shift={(10.4,0)}]
    \draw[acento, line width=1.3pt, line join=round] (0,0) -- (0.9,0) -- (1.05,0.2) -- (1.22,-0.7) -- (1.42,0.85)
      .. controls (1.6,0.35) and (1.8,0.22) .. (2.1,0.24) .. controls (2.5,0.27) and (2.7,0.6) .. (3.0,0.58)
      .. controls (3.3,0.55) and (3.5,0) .. (3.85,0) -- (4.4,0);
    \draw[gristexto, -{Latex[length=1.6mm]}, line width=0.5pt] (2.45,-0.72) node[right, font=\sffamily\scriptsize, text=gristexto] {ST cóncavo elevado} -- (2.07,0.17);
  \end{scope}
\end{tikzpicture}
\caption[Patrones de Brugada en V1--V2]{Representación esquemática de los patrones de Brugada en V1--V2 frente a un complejo normal de V1. El patrón tipo 1 combina una elevación del punto J de al menos \SI{0.2}{\milli\volt}, un ST descendente y convexo y una onda T negativa. El tipo 2 presenta una r' y un ST cóncavo en silla de montar con T positiva. Adaptado de los criterios de consenso \cite{bayesdeluna2012,priori2013}.}
\label{fig:patrones}
\end{figure}

\subsection{Diagnóstico}

El consenso HRS/EHRA/APHRS de 2013 establece que el SBr se diagnostica cuando se observa un patrón tipo~1 en al menos una derivación precordial derecha (V1 o V2, en posición estándar o elevada hasta el segundo espacio intercostal), ya sea espontáneo o inducido por un bloqueante de los canales de sodio como la ajmalina o la flecainida \cite{priori2013}. El consenso de 2016 propuso además una puntuación (\eng{Shanghai score}) que combina el ECG, la clínica, los antecedentes familiares y la genética \cite{antzelevitch2016}. Las guías europeas de 2022 restringieron el valor del patrón inducido por fármacos, que por sí solo ya no basta para el diagnóstico \cite{zeppenfeld2022}.

\section{Implicaciones para el análisis de datos}

\begin{table}[!b]
\centering
\caption[Traducción de los criterios clínicos a variables]{Traducción de los criterios clínicos del patrón de Brugada a variables calculables sobre el latido medio. La definición precisa de cada variable está en el \cref{ap:catalogo}.}
\label{tab:criterios}
\small
\begin{tabularx}{\linewidth}{@{}p{4.2cm}Lp{3.0cm}@{}}
\toprule
\textbf{Criterio clínico} & \textbf{Variable en este trabajo} & \textbf{Nombre} \\
\midrule
Elevación del punto J $\geq$ \SI{0.2}{\milli\volt} & Amplitud del latido en el punto J y a \SIlist{40;80}{\milli\second} & \var{st_j}, \var{st_j40}, \var{st_j80} \\
ST descendente y convexo (\eng{coved}); índice de Corrado & Pendiente del ST entre J y J+\SI{80}{\milli\second}; pendiente final del QRS & \var{st_slope}, \var{qrs_end_slope} \\
Onda T negativa & Mínimo de la onda T e indicador de negatividad & \var{t_neg}, \var{t_negative} \\
Onda r' prominente & Máximo positivo tras la onda S & \var{rprime_amp} \\
Ensanchamiento del QRS & Duración del QRS a partir de la energía de la derivada & \var{qrs_dur_ms} \\
Elevación relativa del ST & Cociente entre el ST y la amplitud del QRS & \var{st_r_ratio} \\
\bottomrule
\end{tabularx}
\end{table}

Las características clínicas del síndrome tienen consecuencias directas sobre cómo hay que preparar los datos:

\begin{enumerate}
  \item \textbf{El patrón es local en el espacio y en el tiempo.} Está en dos derivaciones (V1--V2) y ocupa una ventana de unos \SI{200}{\milli\second} alrededor del punto J. La segmentación y el alineamiento de latidos deben ser precisos en esa zona.
  \item \textbf{El patrón es intermitente y a menudo incompleto.} La mayoría de los pacientes del conjunto de datos no muestra el tipo 1 en su ECG basal, así que las variables deben captar tendencias sutiles y no solo comprobar si se cumple un umbral clínico.
  \item \textbf{Los criterios clínicos son medibles.} Cada elemento del patrón se puede traducir en una variable calculable sobre el latido (\cref{tab:criterios}). Esto permite construir características interpretables y comparar un sistema de reglas explícitas con modelos aprendidos.
\end{enumerate}

